\metatitle{A gamma interpretation for OEIS A046739}
\metadescription{A Foata--Strehl orbit proof for the gamma coefficients of derangement excedance polynomials, using ascending runs and double excedances.}


\section[oeisProofA046739]{A gamma interpretation for A046739}

This page is part of the \hyperref[oeisProofs]{OEIS proof collection}.

Let $\mathcal{D}_n$ be the set of \defin{derangements} of
$[n]=\{1,\dotsc,n\}$, that is, permutations with no fixed point.
An \defin{excedance} of $\sigma$ is an index $i$ with $i<\sigma(i)$;
we write $\operatorname{exc}(\sigma)$ for their number.
The \defin{derangement excedance polynomial}
\[
 D_n(t)\coloneqq\sum_{\sigma\in\mathcal{D}_n}
 t^{\operatorname{exc}(\sigma)}
\]
is the row polynomial of \oeis{A046739}.
An \defin{ascending run} is a maximal consecutive increasing subword
in one-line notation; a singleton run has length one.
A value $i$ is a \defin{double excedance} of $\sigma$ if
$\sigma^{-1}(i)<i<\sigma(i)$.

\begin{theorem}
For $n\geq 1$, we have
\[
 D_n(t)=\sum_{r=0}^{\lfloor n/2\rfloor}
 \xi(n,r)t^r(1+t)^{n-2r},
\]
where $\xi(n,r)$ counts either of the following classes:
\begin{itemize}
\item permutations of $[n]$ with $r$ ascending runs and no singleton run;
\item derangements of $[n]$ with $r$ excedances and no double excedance.
\end{itemize}
\end{theorem}

This is \cite[Thm.~1.4]{AthanasiadisSavvidou2012}.
We outline its orbit proof below.

\subsection[oeisA046739OrbitProof]{The orbit proof}

\begin{proof}
A \defin{descent} of a word $w=w_1\dotsm w_n$ is a position $i<n$
with $w_i>w_{i+1}$; let $\operatorname{des}(w)$ count these positions.
Let $E_n$ consist of permutations in which every left-to-right maximum
(an entry larger than all preceding entries) is followed by a smaller entry.
Write each cycle of a permutation with its largest entry first, order the
cycles by increasing largest entry, and erase the parentheses; call this
fundamental transformation $\varphi$.
Then $\sigma\mapsto\varphi(\sigma^{-1})$ is a bijection
$\mathcal{D}_n\to E_n$ satisfying
$\operatorname{exc}(\sigma)=\operatorname{des}(\varphi(\sigma^{-1}))$.

Set $w_0=0$ and $w_{n+1}=n+1$.
A \defin{double ascent} is an entry $w_i$ with
$w_{i-1}<w_i<w_{i+1}$; a \defin{double descent} has both inequalities reversed.
The modified Foata--Strehl action on $E_n$ moves a double-ascent entry left
to the nearest gap $w_j>w_i>w_{j+1}$, and moves a double-descent entry right
to the nearest gap $w_j<w_i<w_{j+1}$.
These inverse moves preserve $E_n$ and toggle the status of the moved entry.
The toggles indexed by entries commute, and each orbit has a unique
representative with no double descent.

A \defin{peak} satisfies $w_{i-1}<w_i>w_{i+1}$, and a \defin{valley}
satisfies $w_{i-1}>w_i<w_{i+1}$.
If that representative has $r$ descents, it has $r$ peaks and $r$ valleys;
the remaining $n-2r$ entries are double ascents.
Toggling any subset of them increases the descent number by the subset size.
Thus its orbit contributes
\[
 \sum_{u\in\mathcal{O}(w)}t^{\operatorname{des}(u)}
 =t^r(1+t)^{n-2r}.
\]
Under $\varphi(\sigma^{-1})$, the representatives correspond exactly to
derangements with $r$ excedances and no double excedance.
For such a derangement $\sigma$, the word $\varphi(\sigma)$ has $r$
descending runs and no singleton run. Complementing its entries
($i\mapsto n+1-i$) gives the ascending-run model.
Summing the orbit contributions proves the expansion and both interpretations;
see \cite[proof of Thm.~1.4, pp.~18--19]{AthanasiadisSavvidou2012}.
\end{proof}

H.~Sun and Y.~Wang give the same orbit factor directly on derangements
using a cyclic modified Foata--Strehl action
\cite[Thm.~5 and Cor.~6]{SunWang2014}.
Its unique representative has no cyclic double ascent, which is precisely
the condition of having no double excedance defined above.

\subsection[oeisA046739Example]{Example and shifted normalization}

For $n=4$,
\[
 D_4(t)=t+7t^2+t^3=t(1+t)^2+5t^2.
\]
The ascending-run model consists of $1234$ with one run, and
$1324,1423,2314,2413,3412$ with two runs of length two.
Hence $\xi(4,1)=1$ and $\xi(4,2)=5$.

For the shifted normalization $P_m(t)\coloneqq D_{m+4}(t)/t$, $m\geq0$,
write
\[
 P_m(t)=\sum_{j=0}^{\lfloor(m+2)/2\rfloor}
 \gamma(m,j)t^j(1+t)^{m+2-2j}.
\]
Dividing the expansion of $D_{m+4}$ by $t$ and setting $r=j+1$ gives
$\gamma(m,j)=\xi(m+4,j+1)$.
Finally, summing $\xi(n,r)$ over $r$ forgets the number of runs and gives
\oeis{A097899}$(n)$, the number of permutations with no singleton ascending run.
